
\documentclass[11pt,a4paper]{article}

\usepackage{fontspec}
\usepackage[french]{babel}
\usepackage{amsmath,amssymb,amsthm,mathtools}
\usepackage{geometry}
\usepackage{xcolor}
\usepackage{fancyhdr}
\usepackage{enumitem}
\usepackage{hyperref}
\usepackage{titlesec}
\usepackage{microtype}
\usepackage{booktabs}
\usepackage{array}

\geometry{margin=2.2cm}
\setmainfont{DejaVu Serif}
\setsansfont{DejaVu Sans}
\setmonofont{DejaVu Sans Mono}

\definecolor{accent}{HTML}{1F4E79}
\definecolor{lightaccent}{HTML}{EAF2F8}

\pagestyle{fancy}
\fancyhf{}
\lhead{\small Cours Deep Reinforcement Learning}
\chead{\small Classe : 4IA}
\rhead{\small Année : 2026}
\lfoot{\small Enseignant : Mourad Zéraï, PhD - ESPRIT School of Engineering}
\cfoot{}
\rfoot{\thepage}
\setlength{\headheight}{14pt}

\titleformat{\section}{\large\bfseries\color{accent}}{\thesection.}{0.5em}{}
\titleformat{\subsection}{\normalsize\bfseries\color{accent}}{\thesubsection.}{0.5em}{}
\titleformat{\subsubsection}{\normalsize\bfseries}{\thesubsubsection.}{0.5em}{}

\newtheorem*{definition}{Définition}
\newtheorem*{rappel}{Rappel}
\newtheorem*{hypothese}{Hypothèse}
\newtheorem*{propriete}{Propriété}
\newtheorem*{theoreme}{Théorème}
\newtheorem*{remarque}{Remarque}

\newcommand{\E}{\mathbb{E}}
\newcommand{\Pbb}{\mathbb{P}}
\newcommand{\R}{\mathbb{R}}
\newcommand{\argmax}{\operatorname*{argmax}}
\newcommand{\given}{\,|\,}

\begin{document}

\begin{center}
    {\LARGE\bfseries Dérivation et explication pas à pas de Monte Carlo Control}\\[0.4em]
    {\large Note de cours rigoureuse, détaillée et appliquée à FrozenLake}
\end{center}

\vspace{0.6em}

\section{Objectif}

Le but de cette note est d'expliquer \emph{Monte Carlo Control}, c'est-à-dire l'apprentissage d'une politique par estimation d'épisodes complets, sans modèle explicite des transitions, puis de détailler son application à FrozenLake.

\section{Idée de base}

Les méthodes Monte Carlo n'utilisent pas l'équation de Bellman comme mise à jour locale pas à pas. Elles attendent la fin d'un épisode, calculent le retour observé, puis utilisent ce retour comme échantillon de la valeur recherchée.

\begin{definition}[Retour]
\[
G_t = R_{t+1}+\gamma R_{t+2}+\gamma^2 R_{t+3}+\cdots+\gamma^{T-t-1}R_T,
\]
où $T$ désigne l'instant terminal de l'épisode.
\end{definition}

\begin{definition}[Valeur état-action]
\[
q_{\pi}(s,a)=\E_{\pi}[G_t \given S_t=s,A_t=a].
\]
\end{definition}

L'idée est donc simple : si l'on observe beaucoup de retours $G_t$ après être passé par $(s,a)$, leur moyenne devrait approcher $q_{\pi}(s,a)$.

\section{Pourquoi la moyenne empirique est légitime}

Soient
\[
G_t^{(1)},G_t^{(2)},\dots,G_t^{(N)}
\]
des réalisations du retour obtenues lors de visites répétées du couple $(s,a)$ sous la même politique $\pi$.

\begin{theoreme}[Loi forte des grands nombres, version intuitive]
Si les variables sont intégrables et identiquement distribuées, alors
\[
\frac1N\sum_{i=1}^N G_t^{(i)} \to \E_{\pi}[G_t\given S_t=s,A_t=a]
\]
presque sûrement lorsque $N\to\infty$.
\end{theoreme}

Donc un estimateur naturel est
\[
\hat q_N(s,a)=\frac1N\sum_{i=1}^N G_t^{(i)}.
\]

\section{De l'évaluation au contrôle}

\subsection{Évaluation Monte Carlo}

Si la politique $\pi$ est fixée, on peut estimer $q_{\pi}(s,a)$ par moyenne empirique des retours observés après chaque visite de $(s,a)$.

\subsection{Contrôle}

Le contrôle consiste à améliorer la politique au fur et à mesure des estimations. L'idée est :
\begin{enumerate}[leftmargin=1.6em]
    \item estimer $q_{\pi}$ à partir des épisodes ;
    \item rendre la politique plus gloutonne par rapport à cette estimation.
\end{enumerate}

\section{Pourquoi il faut de l'exploration}

Si l'on choisit toujours l'action actuellement jugée meilleure, certains couples $(s,a)$ risquent de ne jamais être revisités. Dans ce cas, leurs estimations ne progresseront plus. Il faut donc garantir un minimum d'exploration.

Deux approches classiques existent :
\begin{itemize}[leftmargin=1.6em]
    \item \emph{exploring starts} : on suppose que tout couple $(s,a)$ peut initier un épisode avec probabilité positive ;
    \item politique \emph{$\varepsilon$-gloutonne} : avec probabilité $\varepsilon$, on choisit une action aléatoire.
\end{itemize}

Dans FrozenLake, la seconde approche est la plus simple à implémenter.

\section{Mise à jour par moyenne incrémentale}

Stocker toutes les visites est possible, mais inutile. Si $Q_n(s,a)$ est la moyenne des $n$ premiers retours et si $G_{n+1}$ est le nouveau retour, alors
\[
Q_{n+1}(s,a)=Q_n(s,a)+\frac{1}{n+1}\left(G_{n+1}-Q_n(s,a)\right).
\]

\paragraph{Justification.}
Par définition,
\[
Q_{n+1}=\frac{nQ_n+G_{n+1}}{n+1}
=Q_n+\frac{G_{n+1}-Q_n}{n+1}.
\]
Cette écriture est une simple réorganisation algébrique.

\section{First-visit Monte Carlo control}

\subsection{Principe}

Dans la variante \emph{first-visit}, pour un épisode donné, on ne met à jour $(s,a)$ qu'à sa première apparition dans l'épisode. Cela évite la corrélation forte entre plusieurs occurrences du même couple dans un même épisode.

\subsection{Étapes}

Pour chaque épisode :
\begin{enumerate}[leftmargin=1.6em]
    \item générer un épisode en suivant la politique courante ;
    \item calculer les retours $G_t$ en remontant l'épisode de la fin vers le début ;
    \item pour chaque premier couple $(s,a)$ visité, mettre à jour $Q(s,a)$ ;
    \item rendre la politique $\varepsilon$-gloutonne par rapport à $Q$.
\end{enumerate}

\section{Pourquoi l'amélioration est naturelle}

Si l'estimateur $Q(s,a)$ approche $q_{\pi}(s,a)$, alors choisir
\[
\pi'(s)\in \argmax_a Q(s,a)
\]
revient à approcher la politique gloutonne par rapport à $q_{\pi}$. On transpose donc l'idée du théorème d'amélioration de politique, mais avec une estimation empirique au lieu d'une quantité exacte.

\section{Application à FrozenLake}

\subsection{Épisode typique}

Un épisode est une suite :
\[
(s_0,a_0,r_1,s_1,a_1,r_2,\dots,s_T).
\]
Dans FrozenLake, la plupart des récompenses valent $0$, sauf l'atteinte du but qui donne généralement $1$.

\subsection{Calcul concret des retours}

Si la fin d'épisode vaut
\[
r_{t+1}=0,\quad r_{t+2}=0,\quad r_{t+3}=1,
\]
alors
\[
G_t = 0 + \gamma\cdot 0 + \gamma^2 \cdot 1 = \gamma^2.
\]

\subsection{Effet du mode glissant}

En mode glissant, deux épisodes partis du même état et de la même action peuvent conduire à des suites très différentes. Monte Carlo est adapté à ce contexte car il apprend directement à partir des trajectoires réalisées, sans exiger le modèle exact.

\subsection{Choix d'une politique $\varepsilon$-gloutonne}

Pour chaque état $s$ :
\[
\pi(a\given s)=
\begin{cases}
1-\varepsilon+\dfrac{\varepsilon}{|\mathcal A|} & \text{si } a \in \argmax_b Q(s,b),\\[0.8em]
\dfrac{\varepsilon}{|\mathcal A|} & \text{sinon.}
\end{cases}
\]
On conserve ainsi une probabilité strictement positive de tester toutes les actions.

\section{Pseudo-code pédagogique}

\begin{verbatim}
Initialiser Q(s,a) arbitrairement
Initialiser la politique comme epsilon-gloutonne par rapport à Q
Initialiser N(s,a)=0

Pour chaque épisode :
    Générer l'épisode sous la politique courante
    G = 0
    Parcourir l'épisode à rebours :
        G = gamma * G + reward
        Si (s,a) est une première visite dans l'épisode :
            N(s,a) = N(s,a) + 1
            Q(s,a) = Q(s,a) + (1/N(s,a)) * (G - Q(s,a))
            Mettre à jour la politique epsilon-gloutonne dans l'état s
\end{verbatim}

\section{Limites et intérêts}

\begin{itemize}[leftmargin=1.6em]
    \item La méthode ne s'applique qu'à des épisodes terminaux ou à des retours bien définis.
    \item Elle n'a pas besoin du modèle de transition.
    \item Elle peut être lente si la récompense utile n'apparaît qu'à la fin, comme dans FrozenLake.
    \item Elle est simple conceptuellement : on apprend par moyenne de retours observés.
\end{itemize}

\section{Ce qu'il faut retenir}

\begin{itemize}[leftmargin=1.6em]
    \item Monte Carlo Control estime directement $q_{\pi}$ par moyenne empirique de retours.
    \item La justification probabiliste repose sur la loi des grands nombres.
    \item L'amélioration de politique se fait en rendant la politique gloutonne ou $\varepsilon$-gloutonne par rapport à $Q$.
    \item Dans FrozenLake, la méthode apprend à partir d'épisodes complets, avec ou sans glissement.
\end{itemize}

\end{document}
